\Needspace{14\baselineskip}
\section{Introduction}
\label{sec:introduction}

% PRISMA Item 3: Rationale — justification in the context of existing knowledge
% PRISMA Item 4: Objectives — explicit statement of the questions addressed

Nine of the 28 included studies on knowledge persistence in coding
agents report adverse outcomes associated with persistence:
retrieval-noise-driven performance drops, model-dependent overhead,
memory poisoning, and evaluation artefacts that complicate the
interpretation of reported gains.
These reports make it important to distinguish when persisting
knowledge is beneficial and when it is not, a decision that precedes
the comparison among technically sophisticated architectures.
Coding agents based on large language models (LLMs) have achieved
growing adoption in software engineering tasks such as automated issue
resolution, code generation, and repository
maintenance~\citep{du2026memory}.
Systems such as SWE-Agent and OpenHands, together with Agentless,
interact iteratively with the source code; they run tests and propose
fixes without direct human intervention.
Agent capability, however, remains limited by a basic constraint:
many benchmarked agents are effectively stateless across tasks unless
an explicit persistence mechanism is added.

The agent rediscovers project architecture, repeats mistakes already
made in earlier sessions, and ignores design decisions previously
validated.
Cross-session knowledge persistence emerged in response to this
limitation: by recording information from past resolutions, the agent
can reuse it in subsequent tasks.

At least 24 persistence systems have been proposed since 2024, all
published between 2024 and the first quarter of 2026.
Architectures range from experience banks that distil resolution
trajectories to knowledge graphs, AST-guided memory, and version-control
metaphors inspired by Git.
These studies do not converge on a unified taxonomy, and most do not
cite contemporary work on the same problem, so evaluation conditions
remain heterogeneous across systems.

From this heterogeneity arise gaps that limit both research and
practice.
No study in the corpus compares three or more persistence architectures
on the same benchmark with the same model;
any claim of architectural superiority therefore relies on comparisons
across studies with heterogeneous experimental conditions.
Cost underreporting compounds the problem: eight of the 24
architecture studies (33\%) report no cost data, rising to 12 of the 28
included studies (43\%) when contextual sources are counted, and no
study in the corpus constructs a Pareto frontier between accuracy and
computational cost.
Nor is there controlled evidence on the relationship between
architectural complexity and effectiveness, so it remains unknown
whether simple architectures are sufficient or whether sophistication
is justified in specific scenarios.

Existing secondary work covers adjacent aspects without addressing
these gaps; the detailed positioning of this review relative to them is
presented in \autoref{sec:related-work}.

This systematic review answers the following primary question:
\textit{how does existing literature address knowledge persistence in
AI-based coding agents, and what empirical evidence exists about the
effectiveness and cost-efficiency of different persistence
architectures?}
To operationalise this primary question, five secondary questions guide
the synthesis, each motivated by a specific gap in the current state of
the field.

\noindent\textbf{SQ1.} Which persistence architectures have been
proposed, and how are they classified by temporal scope,
representational substrate, and control policy?
Architectural heterogeneity prevents both direct comparison and
informed adoption by teams that must choose among competing
alternatives.
An explicit taxonomy reveals which design choices dominate the field
and which remain unexplored.

\noindent\textbf{SQ2.} Which evaluation methods and benchmarks are
used, and how comparable are the results?
When different studies adopt distinct benchmarks, base models, and
experimental protocols, reported gains are not directly comparable with
each other.
This question checks for methodological convergence and identifies the
obstacles that prevent cross-study synthesis.

\noindent\textbf{SQ3.} Which performance results are reported relative
to no-persistence baselines?
The magnitude of the gains determines whether persistence is a marginal
or substantive improvement; synthesising them exposes the extent to
which the gain depends on specific experimental conditions.

\noindent\textbf{SQ4.} Which cost and efficiency metrics are reported,
and what is the relationship between cost and effectiveness?
Persistence adds overhead in storage and retrieval, and it also
consumes the agent's effective context budget.
Without cost data, any accuracy gain is incomplete for practical
adoption, and the cost-effectiveness relationship remains inaccessible
to the reader.

\noindent\textbf{SQ5.} What evidence exists about the relationship
between architectural complexity and effectiveness, including
degradation?
More sophisticated systems do not always outperform simpler
alternatives, and more memory does not always produce more
effectiveness.
Identifying degradation mechanisms and complexity-versus-gain patterns
informs design decisions in scenarios where sophistication is not
justified.

\begin{table}[htbp]
\centering
\caption{Map of the secondary questions: objective of each question and
location of the answer in the manuscript.}
\label{tab:rq-map}
\small
\begin{tabularx}{\textwidth}{l X l l}
\toprule
\textbf{SQ} & \textbf{Objective} & \textbf{Section} & \textbf{Table/Figure} \\
\midrule
SQ1 & Map and classify the persistence architectures &
  \autoref{sec:sq1} & \autoref{tab:architectures}, \autoref{fig:taxonomy} \\
SQ2 & Catalogue evaluation methods and benchmarks &
  \autoref{sec:sq2} & \autoref{tab:benchmarks} \\
SQ3 & Synthesise performance results vs.\ baselines &
  \autoref{sec:sq3} & \autoref{tab:performance}, \autoref{fig:gains} \\
SQ4 & Quantify coverage and cost-effectiveness &
  \autoref{sec:sq4} & \autoref{tab:cost}, \autoref{fig:cost} \\
SQ5 & Assess the relationship between complexity and effectiveness &
  \autoref{sec:sq5} & \autoref{tab:degradation} \\
\bottomrule
\end{tabularx}
\end{table}

\noindent The scope covers publications from 2023 to March 2026, with
searches in Scopus and arXiv complemented by citation chasing.

The contributions of this review are as follows:

\begin{enumerate}
  \item the first systematic mapping of the field, classifying 24
    persistence architectures for coding agents along three dimensions:
    temporal scope of memory, representational substrate, and
    read/write policy;
  \item a quantitative synthesis of 19 controlled effectiveness
    comparisons (18 positive, one controlled negative result from
    CTIM-Rover), paired with a cost-reporting synthesis of the studies
    that publish monetary, token, or latency data;
    together they identify diminishing returns on strong baselines and
    a bimodal impact on operating cost;
  \item empirical confirmation of three central gaps in the field:
    absence of head-to-head comparisons across architectures on the
    same benchmark, systematic underreporting of per-task cost, and
    insufficient coverage of the relationship between architectural
    complexity and effectiveness.
\end{enumerate}

The confirmed gaps motivate a controlled empirical comparative study,
discussed in \autoref{sec:agenda-pesquisa}.
